% ============================================================
\chapter*{Pr\'eface}
\addcontentsline{toc}{chapter}{Pr\'eface}
\markboth{Pr\'eface}{Pr\'eface}
% ============================================================

\section*{Objectifs du cours}

La \textbf{finance quantitative} se situe au carrefour des math\'ematiques,
de la statistique et de l'\'economie financi\`ere. Ce cours de niveau Master
a pour ambition de fournir aux \'etudiants une formation rigoureuse couvrant
les fondements th\'eoriques et les m\'ethodes num\'eriques indispensables \`a
la compr\'ehension et \`a la pratique de la finance de march\'e moderne.

Les march\'es financiers contemporains sont caract\'eris\'es par une complexit\'e
croissante~: produits d\'eriv\'es exotiques, trading algorithmique \`a haute
fr\'equence, gestion quantitative de portefeuilles, et plus r\'ecemment
l'irruption de l'apprentissage automatique dans la mod\'elisation financi\`ere.
Ce cours vise \`a donner aux \'etudiants les outils math\'ematiques et
informatiques n\'ecessaires pour \'evoluer dans cet environnement.

L'approche adopt\'ee est r\'esolument \textbf{math\'ematique et computationnelle}~:
chaque concept th\'eorique est accompagn\'e d'une impl\'ementation Python
permettant de v\'erifier les r\'esultats analytiques et de d\'evelopper
l'intuition num\'erique.

\section*{Pr\'erequis}

Ce cours suppose que l'\'etudiant poss\`ede des connaissances solides dans
les domaines suivants~:

\begin{itemize}
    \item \textbf{Analyse math\'ematique}~: calcul diff\'erentiel et int\'egral,
          s\'eries, espaces de fonctions, \'equations diff\'erentielles ordinaires.
    \item \textbf{Alg\`ebre lin\'eaire}~: espaces vectoriels, matrices,
          valeurs propres, d\'ecompositions matricielles (Cholesky, SVD).
    \item \textbf{Probabilit\'es}~: espaces de probabilit\'e, variables
          al\'eatoires, lois classiques, esp\'erance conditionnelle,
          th\'eor\`emes limites (loi des grands nombres, TCL).
    \item \textbf{Statistique}~: estimation param\'etrique, tests d'hypoth\`eses,
          r\'egression lin\'eaire, maximum de vraisemblance.
    \item \textbf{Programmation}~: ma\^itrise de Python (ou d'un langage
          scientifique \'equivalent), notions de programmation orient\'ee objet.
\end{itemize}

\section*{Organisation du cours}

Le cours est organis\'e en onze chapitres, chacun con\c{c}u pour \^etre
trait\'e en une \`a deux semaines de cours magistral accompagn\'e de
travaux dirig\'es et de s\'eances de travaux pratiques sur ordinateur.

\begin{description}
    \item[Chapitre 1 --- March\'es financiers et mod\'elisation.]
        Introduction aux march\'es, actifs financiers, hypoth\`eses de
        mod\'elisation, absence d'opportunit\'e d'arbitrage (AOA),
        mod\`eles en temps discret (Cox--Ross--Rubinstein).

    \item[Chapitre 2 --- Mod\`ele de Black--Scholes.]
        D\'erivation du mod\`ele, \'equation aux d\'eriv\'ees partielles
        de Black--Scholes, formules ferm\'ees pour les options europ\'eennes,
        calcul et interpr\'etation des Grecques.

    \item[Chapitre 3 --- M\'ethodes de pricing.]
        Arbres binomiaux et trinomiaux, m\'ethodes de Monte Carlo,
        diff\'erences finies, techniques de r\'eduction de variance.

    \item[Chapitre 4 --- Couverture et gestion des risques.]
        Delta-hedging, couverture dynamique, gamma-hedging, co\^ut
        de la couverture, limites pratiques du mod\`ele de Black--Scholes.

    \item[Chapitre 5 --- Calcul stochastique en finance.]
        Mouvement brownien, int\'egrale d'It\^o, formule d'It\^o,
        th\'eor\`eme de Girsanov, repr\'esentation des martingales.

    \item[Chapitre 6 --- Volatilit\'e stochastique.]
        Smile et surface de volatilit\'e, mod\`ele de Heston, mod\`ele
        SABR, calibration, volatilit\'e locale vs stochastique.

    \item[Chapitre 7 --- Mod\`eles de taux d'int\'er\^et.]
        Structure par terme, mod\`eles de Vasicek, CIR, Hull--White,
        cadre HJM, pricing d'obligations et de caps/floors.

    \item[Chapitre 8 --- Optimisation de portefeuille.]
        Th\'eorie de Markowitz, fronti\`ere efficiente, mod\`ele CAPM,
        mod\`ele de Black--Litterman, optimisation robuste.

    \item[Chapitre 9 --- Mesures de risque.]
        Value-at-Risk (VaR), Expected Shortfall (CVaR), propri\'et\'es
        de coh\'erence, m\'ethodes de calcul, backtesting, stress testing.

    \item[Chapitre 10 --- D\'eriv\'es exotiques.]
        Options barri\`eres, options asiatiques, lookback, options sur
        panier, pricing analytique et num\'erique.

    \item[Chapitre 11 --- Apprentissage automatique en finance.]
        R\'eseaux de neurones pour le pricing, apprentissage par renforcement
        pour le trading, d\'etection d'anomalies, risques de sur-apprentissage.
\end{description}

\section*{M\'ethodologie p\'edagogique}

Chaque chapitre suit une structure coh\'erente~:

\begin{enumerate}
    \item \textbf{Th\'eorie}~: pr\'esentation rigoureuse des concepts
          math\'ematiques avec d\'emonstrations compl\`etes ou esquiss\'ees.
    \item \textbf{Interpr\'etation financi\`ere}~: discussion de la
          signification \'economique des r\'esultats et de leurs limites.
    \item \textbf{Impl\'ementation}~: code Python illustrant les m\'ethodes
          num\'eriques avec les biblioth\`eques \texttt{numpy},
          \texttt{scipy}, \texttt{matplotlib}.
    \item \textbf{Exercices}~: probl\`emes th\'eoriques et num\'eriques
          de difficult\'e progressive pour approfondir la compr\'ehension.
\end{enumerate}

\section*{Conventions de notation}

Tout au long de ce cours, nous utiliserons les notations suivantes~:

\begin{center}
\begin{tabular}{cl}
    \toprule
    \textbf{Notation} & \textbf{Signification} \\
    \midrule
    $(\Omega, \mathcal{F}, \PP)$ & Espace de probabilit\'e \\
    $\E[X]$ & Esp\'erance de la variable al\'eatoire $X$ \\
    $\E^\mathbb{Q}[\cdot]$ & Esp\'erance sous la mesure risque-neutre \\
    $\Var(X)$ & Variance de $X$ \\
    $\Cov(X,Y)$ & Covariance de $X$ et $Y$ \\
    $\R, \R^+$ & Nombres r\'eels, r\'eels positifs \\
    $\N$ & Nombres naturels \\
    $W_t$ ou $B_t$ & Mouvement brownien standard \\
    $(\mathcal{F}_t)_{t \geq 0}$ & Filtration naturelle \\
    $\mathbb{Q}$ & Mesure de probabilit\'e risque-neutre \\
    $r$ & Taux d'int\'er\^et sans risque (continu) \\
    $\sigma$ & Volatilit\'e de l'actif sous-jacent \\
    $S_t$ & Prix de l'actif sous-jacent \`a l'instant $t$ \\
    $C(S,t), P(S,t)$ & Prix du call et du put europ\'eens \\
    $K$ & Prix d'exercice (strike) \\
    $T$ & Maturit\'e de l'option \\
    $\Delta, \Gamma, \Theta, \mathcal{V}, \rho$ & Lettres grecques (sensibilit\'es) \\
    \bottomrule
\end{tabular}
\end{center}

\section*{Logiciels et outils}

Les impl\'ementations num\'eriques de ce cours utilisent l'\'ecosyst\`eme
Python scientifique~:

\begin{itemize}
    \item \texttt{numpy}~: calcul matriciel, alg\`ebre lin\'eaire, g\'en\'eration
          de nombres al\'eatoires.
    \item \texttt{scipy}~: optimisation (\texttt{scipy.optimize}), int\'egration,
          fonctions sp\'eciales, loi normale (\texttt{scipy.stats.norm}).
    \item \texttt{matplotlib}~: visualisation graphique, trac\'e de surfaces
          de volatilit\'e, fronti\`eres efficientes.
    \item \texttt{pandas}~: manipulation et analyse de s\'eries temporelles
          financi\`eres.
    \item \texttt{scikit-learn}~: algorithmes d'apprentissage automatique
          classiques (chapitre~11).
    \item \texttt{tensorflow} / \texttt{pytorch}~: r\'eseaux de neurones
          profonds (chapitre~11).
\end{itemize}

\section*{Th\`emes transversaux}

Plusieurs th\`emes traversent l'ensemble du cours et m\'eritent d'\^etre
signal\'es d\`es cette pr\'eface~:

\begin{itemize}
    \item \textbf{Absence d'arbitrage}~: principe fondamental qui sous-tend
          toute la th\'eorie du pricing en finance quantitative.
    \item \textbf{Compl\'etude des march\'es}~: la question de savoir si
          tout actif contingent peut \^etre r\'epliqu\'e par une strat\'egie
          auto-finan\c{c}ante.
    \item \textbf{Risque de mod\`ele}~: les cons\'equences de l'utilisation
          de mod\`eles math\'ematiques simplifi\'es pour des r\'ealit\'es
          financi\`eres complexes.
    \item \textbf{Calibration vs estimation}~: la distinction entre calibrer
          un mod\`ele sur les prix de march\'e et estimer ses param\`etres
          \`a partir de donn\'ees historiques.
\end{itemize}

\section*{R\'ef\'erences bibliographiques}

Ce cours s'appuie sur les ouvrages de r\'ef\'erence suivants~:

\begin{enumerate}
    \item \textsc{Shreve}, S.E. --- \textit{Stochastic Calculus for Finance
          I \& II}, Springer, 2004.
    \item \textsc{Hull}, J.C. --- \textit{Options, Futures, and Other Derivatives},
          Pearson, 2017 (10\textsuperscript{e} \'edition).
    \item \textsc{Bj\"ork}, T. --- \textit{Arbitrage Theory in Continuous Time},
          Oxford University Press, 2009 (3\textsuperscript{e} \'edition).
    \item \textsc{Glasserman}, P. --- \textit{Monte Carlo Methods in Financial
          Engineering}, Springer, 2003.
    \item \textsc{Gatheral}, J. --- \textit{The Volatility Surface: A Practitioner's
          Guide}, Wiley, 2006.
    \item \textsc{Brigo}, D. \& \textsc{Mercurio}, F. --- \textit{Interest Rate
          Models --- Theory and Practice}, Springer, 2006.
    \item \textsc{McNeil}, A.J., \textsc{Frey}, R. \& \textsc{Embrechts}, P.
          --- \textit{Quantitative Risk Management: Concepts, Techniques and Tools},
          Princeton University Press, 2015.
    \item \textsc{L\'opez de Prado}, M. --- \textit{Advances in Financial Machine
          Learning}, Wiley, 2018.
    \item \textsc{Lamberton}, D. \& \textsc{Lapeyre}, B. --- \textit{Introduction
          au calcul stochastique appliqu\'e \`a la finance}, Ellipses, 2012.
    \item \textsc{Dana}, R.-A. \& \textsc{Jeanblanc}, M. --- \textit{March\'es
          financiers en temps continu}, Economica, 2003.
\end{enumerate}

\section*{Remerciements}

Ce cours a b\'en\'efici\'e des contributions de nombreux coll\`egues
enseignants-chercheurs en math\'ematiques financi\`eres. Nous remercions
particuli\`erement les \'etudiants des promotions pr\'ec\'edentes dont
les questions et remarques pertinentes ont permis d'am\'eliorer
consid\'erablement la pr\'esentation de ce mat\'eriel p\'edagogique.

\bigskip
\begin{flushright}
    \textit{L'auteur}\\
    \textit{Mars 2026}
\end{flushright}

\newpage
